\documentclass[12pt,a4paper]{article}

\usepackage[T2A]{fontenc}
\usepackage[utf8]{inputenc}
\usepackage[russian]{babel}
\usepackage{amsmath,amssymb,amsthm,mathtools}
\usepackage{enumitem}
\usepackage{listings}
\usepackage{xcolor}
\usepackage{geometry}

\geometry{margin=2.2cm}

\newtheorem{corollary}{Следствие}
\newtheorem{definition}{Определение}
\newtheorem{theorem}{Теорема}
\newtheorem{lemma}{Лемма}
\newtheorem{example}{Пример}
\newtheorem{remark}{Замечание}

\lstdefinestyle{fpstyle}{
  basicstyle=\ttfamily\small,
  keywordstyle=\color{blue},
  commentstyle=\color{gray},
  stringstyle=\color{teal},
  showstringspaces=false,
  frame=single,
  breaklines=true,
  tabsize=2,
  columns=fullflexible
  literate=
    {А}{{\CYRA}}1
    {Б}{{\CYRB}}1
    {В}{{\CYRV}}1
    {Г}{{\CYRG}}1
    {Д}{{\CYRD}}1
    {Е}{{\CYRE}}1
    {Ё}{{\CYRYO}}1
    {Ж}{{\CYRZH}}1
    {З}{{\CYRZ}}1
    {И}{{\CYRI}}1
    {Й}{{\CYRISHRT}}1
    {К}{{\CYRK}}1
    {Л}{{\CYRL}}1
    {М}{{\CYRM}}1
    {Н}{{\CYRN}}1
    {О}{{\CYRO}}1
    {П}{{\CYRP}}1
    {Р}{{\CYRR}}1
    {С}{{\CYRS}}1
    {Т}{{\CYRT}}1
    {У}{{\CYRU}}1
    {Ф}{{\CYRF}}1
    {Х}{{\CYRH}}1
    {Ц}{{\CYRC}}1
    {Ч}{{\CYRCH}}1
    {Ш}{{\CYRSH}}1
    {Щ}{{\CYRSHCH}}1
    {Ъ}{{\CYRHRDSN}}1
    {Ы}{{\CYRERY}}1
    {Ь}{{\CYRSFTSN}}1
    {Э}{{\CYREREV}}1
    {Ю}{{\CYRYU}}1
    {Я}{{\CYRYA}}1
    {а}{{\cyra}}1
    {б}{{\cyrb}}1
    {в}{{\cyrv}}1
    {г}{{\cyrg}}1
    {д}{{\cyrd}}1
    {е}{{\cyre}}1
    {ё}{{\cyryo}}1
    {ж}{{\cyrzh}}1
    {з}{{\cyrz}}1
    {и}{{\cyri}}1
    {й}{{\cyrishrt}}1
    {к}{{\cyrk}}1
    {л}{{\cyrl}}1
    {м}{{\cyrm}}1
    {н}{{\cyrn}}1
    {о}{{\cyro}}1
    {п}{{\cyrp}}1
    {р}{{\cyrr}}1
    {с}{{\cyrs}}1
    {т}{{\cyrt}}1
    {у}{{\cyru}}1
    {ф}{{\cyrf}}1
    {х}{{\cyrh}}1
    {ц}{{\cyrc}}1
    {ч}{{\cyrch}}1
    {ш}{{\cyrsh}}1
    {щ}{{\cyrshch}}1
    {ъ}{{\cyrhrdsn}}1
    {ы}{{\cyrery}}1
    {ь}{{\cyrsftsn}}1
    {э}{{\cyrerev}}1
    {ю}{{\cyryu}}1
    {я}{{\cyrya}}1
}

\lstset{style=fpstyle}

\title{Билет №46 \\ \small Функциональное программирование. Функциональная декомпозиция. Персистентные структуры данных. Подходы к проектированию и реализации функциональных программ. (ИТМО)}
\author{Сериков Владислав Александрович}
\date{13 мая 2026}

\begin{document}

\maketitle
\tableofcontents
\newpage

\section{Парадигма функционального программирования}

\subsection{Основная идея}

\begin{definition}
\textbf{Функциональное программирование} --- парадигма программирования, в которой вычисление рассматривается как вычисление значений математических функций, а построение программы основано на композиции функций, рекурсии, неизменяемых данных и минимизации побочных эффектов.
\end{definition}

В императивном программировании программа обычно понимается как последовательность команд, изменяющих состояние памяти. В функциональном программировании программа понимается как выражение, которое вычисляется до значения.

\[
\text{Императивный стиль: } \quad
s := 0;\quad
\text{for } x \in xs \text{ do } s := s + x
\]

\[
\text{Функциональный стиль: } \quad
\operatorname{sum}(xs) =
\begin{cases}
0, & xs = [] \\
x + \operatorname{sum}(xs'), & xs = x : xs'
\end{cases}
\]

Функциональная программа описывает не столько \emph{как} менять состояние, сколько \emph{что} вычислять.

\subsection{Ключевые свойства функционального программирования}

К основным свойствам функционального программирования относятся:

\begin{enumerate}
  \item \textbf{Функции как значения первого класса.}

  Функции можно:
  \begin{itemize}
    \item передавать как аргументы;
    \item возвращать как результат;
    \item сохранять в структурах данных;
    \item создавать динамически.
  \end{itemize}

  \item \textbf{Функции высшего порядка.}

  \begin{definition}
  \textbf{Функция высшего порядка} --- функция, которая принимает другие функции в качестве аргументов или возвращает функцию в качестве результата.
  \end{definition}

  Примеры: \texttt{map}, \texttt{filter}, \texttt{fold}.

  \item \textbf{Чистые функции.}

  \begin{definition}
  Функция называется \textbf{чистой}, если:
  \begin{enumerate}
    \item её результат зависит только от значений аргументов;
    \item она не имеет наблюдаемых побочных эффектов.
  \end{enumerate}
  \end{definition}

  Пример чистой функции:
  \[
  f(x) = x^2 + 1.
  \]

  Пример нечистой функции:
  \begin{lstlisting}[language=Python]
counter = 0

def next_id():
    global counter
    counter += 1
    return counter
  \end{lstlisting}

  \item \textbf{Ссылочная прозрачность.}

  \begin{definition}
  Выражение называется \textbf{ссылочно прозрачным}, если его можно заменить его значением без изменения поведения программы.
  \end{definition}

  Например, если
  \[
  f(2) = 5,
  \]
  то в чистом функциональном языке выражение \texttt{f(2)} можно заменить на \texttt{5}.

  \item \textbf{Неизменяемость данных.}

  После создания значение не изменяется. Операции обновления создают новое значение, возможно разделяющее часть структуры со старым.

  \item \textbf{Рекурсия вместо циклов.}

  Поскольку переменные обычно не изменяются, повторение выражается через рекурсию или функции высшего порядка.

  \item \textbf{Композиция функций.}

  Сложные вычисления строятся из простых функций.

  \[
  (g \circ f)(x) = g(f(x)).
  \]
\end{enumerate}

\subsection{Чистое и нечистое функциональное программирование}

Функциональные языки можно условно разделить на две группы.

\begin{enumerate}
  \item \textbf{Чистые функциональные языки.}

  В них побочные эффекты строго контролируются системой типов или специальными механизмами.

  Примеры:
  \begin{itemize}
    \item Haskell;
    \item Clean;
    \item Idris.
  \end{itemize}

  \item \textbf{Нечистые или мультипарадигменные функциональные языки.}

  В них функциональный стиль поддерживается, но побочные эффекты разрешены напрямую.

  Примеры:
  \begin{itemize}
    \item OCaml;
    \item F\#;
    \item Scala;
    \item Clojure;
    \item Lisp;
    \item Scheme;
    \item JavaScript;
    \item Python.
  \end{itemize}
\end{enumerate}

\subsection{Преимущества функционального подхода}

Функциональный подход полезен по следующим причинам:

\begin{enumerate}
  \item \textbf{Локальность рассуждений.}

  Чистую функцию можно анализировать отдельно, не учитывая глобальное состояние.

  \item \textbf{Упрощение тестирования.}

  Для чистой функции достаточно проверить соответствие входов и выходов.

  \item \textbf{Упрощение параллелизма.}

  Если данные неизменяемы и нет побочных эффектов, отсутствуют гонки за изменяемое состояние.

  \item \textbf{Модульность.}

  Функции высшего порядка и ленивые вычисления позволяют разделять программу на независимые компоненты.

  \item \textbf{Повышенная надёжность.}

  Меньше неявного состояния --- меньше ошибок, связанных с порядком выполнения.
\end{enumerate}

\subsection{Недостатки и ограничения}

К сложностям функционального программирования относятся:

\begin{enumerate}
  \item возможные накладные расходы на создание новых структур вместо изменения старых;
  \item необходимость понимать рекурсию, замыкания, ленивость, монады и другие абстракции;
  \item сложность взаимодействия с внешним миром, поскольку ввод-вывод по природе имеет побочные эффекты;
  \item иногда менее предсказуемая производительность из-за ленивости и разделения структур;
  \item необходимость оптимизаций компилятора: устранение хвостовой рекурсии, размораживание ленивых вычислений, инлайнинг.
\end{enumerate}

\section{Математические основы функционального программирования}

\subsection{Лямбда-исчисление}

\begin{definition}
\textbf{Лямбда-исчисление} --- формальная система для описания функций, их применения и вычисления.
\end{definition}

Синтаксис нетипизированного лямбда-исчисления:

\[
M ::= x \mid \lambda x.M \mid M\,N.
\]

Здесь:
\begin{itemize}
  \item $x$ --- переменная;
  \item $\lambda x.M$ --- абстракция, то есть функция с параметром $x$ и телом $M$;
  \item $M\,N$ --- применение функции $M$ к аргументу $N$.
\end{itemize}

Пример:

\[
(\lambda x. x + 1)\ 5 \longrightarrow 6.
\]

\subsection{Свободные и связанные переменные}

\begin{definition}
В выражении $\lambda x.M$ переменная $x$ является \textbf{связанной} в $M$. Переменная, не связанная никакой лямбда-абстракцией, называется \textbf{свободной}.
\end{definition}

Пример:

\[
\lambda x. x + y.
\]

Здесь $x$ --- связанная переменная, а $y$ --- свободная.

\subsection{Подстановка}

Подстановка $M[x := N]$ означает замену свободных вхождений $x$ в $M$ на выражение $N$.

Например:

\[
(x + 1)[x := 5] = 5 + 1.
\]

Однако при подстановке необходимо избегать захвата переменных.

Пример опасной подстановки:

\[
(\lambda y. x)[x := y].
\]

Если выполнить её напрямую, получится:

\[
\lambda y. y,
\]

но свободная переменная $y$ из подставляемого выражения стала связанной. Поэтому перед подстановкой требуется переименование связанной переменной:

\[
\lambda y. x \equiv \lambda z. x,
\]

и тогда:

\[
(\lambda z. x)[x := y] = \lambda z. y.
\]

\subsection{Альфа-, бета- и эта-преобразования}

\subsubsection{Альфа-преобразование}

\begin{definition}
\textbf{Альфа-преобразование} --- переименование связанной переменной.
\end{definition}

\[
\lambda x. x \equiv_\alpha \lambda y. y.
\]

\subsubsection{Бета-редукция}

\begin{definition}
\textbf{Бета-редукция} --- правило применения функции к аргументу:
\[
(\lambda x.M)\ N \to_\beta M[x := N].
\]
\end{definition}

Пример:

\[
(\lambda x. x^2 + 1)\ 3
\to_\beta 3^2 + 1
= 10.
\]

\subsubsection{Эта-преобразование}

\begin{definition}
\textbf{Эта-преобразование} выражает экстенсиональность функций:
\[
\lambda x. f\,x \equiv_\eta f,
\]
если $x$ не входит свободно в $f$.
\end{definition}

Идея: функция, которая только применяет $f$ к аргументу, эквивалентна самой $f$.

\subsection{Нормальная форма}

\begin{definition}
Лямбда-терм находится в \textbf{нормальной форме}, если к нему нельзя применить ни одну бета-редукцию.
\end{definition}

Пример нормальной формы:

\[
\lambda x. x.
\]

Пример терма, не находящегося в нормальной форме:

\[
(\lambda x. x)\ y.
\]

Он редуцируется:

\[
(\lambda x. x)\ y \to y.
\]

\subsection{Теорема Чёрча--Россера}

\begin{theorem}[Чёрч--Россер]
Если лямбда-терм $M$ редуцируется к термам $N_1$ и $N_2$, то существует терм $P$, такой что $N_1$ и $N_2$ редуцируются к $P$:
\[
M \to_\beta^* N_1,\quad M \to_\beta^* N_2
\Rightarrow
\exists P:\ N_1 \to_\beta^* P,\quad N_2 \to_\beta^* P.
\]
\end{theorem}

\begin{proof}
Приведём доказательство через свойство ромба для параллельной редукции.

Введём отношение параллельной редукции $\Rightarrow$, которое за один шаг редуцирует произвольное множество непересекающихся бета-редексов.

Оно определяется индуктивно:
\begin{enumerate}
  \item $x \Rightarrow x$;
  \item если $M \Rightarrow M'$, то $\lambda x.M \Rightarrow \lambda x.M'$;
  \item если $M \Rightarrow M'$ и $N \Rightarrow N'$, то $M\,N \Rightarrow M'\,N'$;
  \item если $M \Rightarrow M'$ и $N \Rightarrow N'$, то
  \[
  (\lambda x.M)\,N \Rightarrow M'[x := N'].
  \]
\end{enumerate}

Докажем, что параллельная редукция обладает свойством ромба:
\[
M \Rightarrow M_1,\quad M \Rightarrow M_2
\Rightarrow
\exists P:\ M_1 \Rightarrow P,\quad M_2 \Rightarrow P.
\]

Доказательство проводится структурной индукцией по терму $M$.

\textbf{Случай 1.} $M = x$.

Тогда единственно возможно:
\[
x \Rightarrow x.
\]
Следовательно, $M_1 = M_2 = x$, и можно взять $P = x$.

\textbf{Случай 2.} $M = \lambda x.A$.

Тогда:
\[
M_1 = \lambda x.A_1,\quad M_2 = \lambda x.A_2,
\]
где
\[
A \Rightarrow A_1,\quad A \Rightarrow A_2.
\]
По предположению индукции существует $B$, такой что:
\[
A_1 \Rightarrow B,\quad A_2 \Rightarrow B.
\]
Тогда:
\[
\lambda x.A_1 \Rightarrow \lambda x.B,\quad
\lambda x.A_2 \Rightarrow \lambda x.B.
\]
Берём $P = \lambda x.B$.

\textbf{Случай 3.} $M = A\,B$, причём $A$ не обязательно является абстракцией.

Если верхний редекс не сокращается, то:
\[
M_1 = A_1 B_1,\quad M_2 = A_2 B_2,
\]
где
\[
A \Rightarrow A_1,\quad A \Rightarrow A_2,\quad
B \Rightarrow B_1,\quad B \Rightarrow B_2.
\]
По предположению индукции существуют $C$ и $D$, такие что:
\[
A_1 \Rightarrow C,\quad A_2 \Rightarrow C,
\]
\[
B_1 \Rightarrow D,\quad B_2 \Rightarrow D.
\]
Тогда:
\[
A_1B_1 \Rightarrow CD,\quad A_2B_2 \Rightarrow CD.
\]
Берём $P = CD$.

\textbf{Случай 4.} $M = (\lambda x.A)B$.

Здесь возможны два типа параллельной редукции:
\begin{enumerate}
  \item редуцируются только части;
  \item редуцируется внешний бета-редекс.
\end{enumerate}

Пусть:
\[
(\lambda x.A)B \Rightarrow (\lambda x.A_1)B_1
\]
и
\[
(\lambda x.A)B \Rightarrow A_2[x := B_2],
\]
где
\[
A \Rightarrow A_1,\quad A \Rightarrow A_2,\quad
B \Rightarrow B_1,\quad B \Rightarrow B_2.
\]

По предположению индукции существуют $C$ и $D$:
\[
A_1 \Rightarrow C,\quad A_2 \Rightarrow C,
\]
\[
B_1 \Rightarrow D,\quad B_2 \Rightarrow D.
\]

Тогда:
\[
(\lambda x.A_1)B_1 \Rightarrow C[x := D],
\]
и также:
\[
A_2[x := B_2] \Rightarrow C[x := D].
\]

Последний переход следует из леммы о совместимости параллельной редукции с подстановкой:
если $A_2 \Rightarrow C$ и $B_2 \Rightarrow D$, то
\[
A_2[x := B_2] \Rightarrow C[x := D].
\]

Следовательно, для всех случаев существует общий потомок $P$.

Теперь заметим, что обычная бета-редукция $\to_\beta$ является частным случаем параллельной редукции, а любая параллельная редукция реализуется конечной последовательностью обычных бета-редукций:
\[
\to_\beta \subseteq \Rightarrow \subseteq \to_\beta^*.
\]

Из свойства ромба для $\Rightarrow$ следует свойство слияния для $\to_\beta^*$:
если
\[
M \to_\beta^* N_1,\quad M \to_\beta^* N_2,
\]
то существует $P$:
\[
N_1 \to_\beta^* P,\quad N_2 \to_\beta^* P.
\]

Это и есть утверждение теоремы Чёрча--Россера.
\end{proof}

\begin{corollary}
Если лямбда-терм имеет нормальную форму, то она единственна с точностью до альфа-переименования.
\end{corollary}

\begin{proof}
Пусть терм $M$ имеет две нормальные формы $N_1$ и $N_2$:
\[
M \to_\beta^* N_1,\quad M \to_\beta^* N_2.
\]

По теореме Чёрча--Россера существует терм $P$, такой что:
\[
N_1 \to_\beta^* P,\quad N_2 \to_\beta^* P.
\]

Но $N_1$ и $N_2$ нормальны, поэтому из них нельзя выполнить бета-редукцию. Следовательно:
\[
P = N_1 = N_2
\]
с точностью до альфа-переименования.
\end{proof}

\section{Функции высшего порядка}

\subsection{Отображение: map}

\begin{definition}
Функция \texttt{map} применяет функцию к каждому элементу списка.
\end{definition}

Математически:

\[
\operatorname{map}\ f\ [x_1,x_2,\dots,x_n]
=
[f(x_1), f(x_2), \dots, f(x_n)].
\]

Рекурсивное определение:

\[
\operatorname{map}\ f\ [] = [],
\]

\[
\operatorname{map}\ f\ (x:xs) = f(x) : \operatorname{map}\ f\ xs.
\]

Минимальный пример:

\[
\operatorname{map}\ (x \mapsto x^2)\ [1,2,3] = [1,4,9].
\]

Алгоритм \texttt{map}:

\begin{enumerate}
  \item Если список пуст, вернуть пустой список.
  \item Иначе взять голову списка $x$.
  \item Применить функцию $f$ к $x$.
  \item Рекурсивно обработать хвост списка.
  \item Построить новый список из $f(x)$ и результата рекурсии.
\end{enumerate}

Иллюстрация:

\[
[1,2,3]
\overset{x^2}{\longmapsto}
[1^2,2^2,3^2]
=
[1,4,9].
\]

Пример на Haskell-подобном псевдокоде:

\begin{lstlisting}
map f []     = []
map f (x:xs) = f x : map f xs
\end{lstlisting}

\subsection{Фильтрация: filter}

\begin{definition}
Функция \texttt{filter} оставляет в списке только элементы, удовлетворяющие предикату.
\end{definition}

\[
\operatorname{filter}\ p\ [] = [],
\]

\[
\operatorname{filter}\ p\ (x:xs) =
\begin{cases}
x : \operatorname{filter}\ p\ xs, & p(x) = \text{true}, \\
\operatorname{filter}\ p\ xs, & p(x) = \text{false}.
\end{cases}
\]

Пример:

\[
\operatorname{filter}\ \operatorname{even}\ [1,2,3,4,5,6]
=
[2,4,6].
\]

Алгоритм \texttt{filter}:

\begin{enumerate}
  \item Если список пуст, вернуть пустой список.
  \item Иначе взять первый элемент $x$.
  \item Проверить условие $p(x)$.
  \item Если условие истинно, включить $x$ в результат.
  \item Если условие ложно, пропустить $x$.
  \item Рекурсивно обработать хвост списка.
\end{enumerate}

\subsection{Свёртка: fold}

\begin{definition}
\textbf{Свёртка} --- функция высшего порядка, обобщающая рекурсивную обработку структуры данных.
\end{definition}

Правая свёртка:

\[
\operatorname{foldr}\ f\ z\ [x_1,x_2,\dots,x_n]
=
f(x_1, f(x_2, \dots f(x_n,z)\dots)).
\]

Рекурсивно:

\[
\operatorname{foldr}\ f\ z\ [] = z,
\]

\[
\operatorname{foldr}\ f\ z\ (x:xs) = f(x, \operatorname{foldr}\ f\ z\ xs).
\]

Левая свёртка:

\[
\operatorname{foldl}\ f\ z\ [x_1,x_2,\dots,x_n]
=
f(\dots f(f(z,x_1),x_2)\dots,x_n).
\]

Пример суммы:

\[
\operatorname{foldr}\ (+)\ 0\ [1,2,3]
=
1 + (2 + (3 + 0))
=
6.
\]

Пример произведения:

\[
\operatorname{foldr}\ (*)\ 1\ [2,3,4]
=
2 \cdot (3 \cdot (4 \cdot 1))
=
24.
\]

Алгоритм \texttt{foldr}:

\begin{enumerate}
  \item Если список пуст, вернуть начальное значение $z$.
  \item Иначе взять голову $x$.
  \item Рекурсивно свернуть хвост списка.
  \item Применить функцию $f$ к $x$ и результату свёртки хвоста.
\end{enumerate}

\subsection{Закон композиции для map}

\begin{theorem}
Для любых функций $f$ и $g$ и любого списка $xs$ выполняется:
\[
\operatorname{map}\ f\ (\operatorname{map}\ g\ xs)
=
\operatorname{map}\ (f \circ g)\ xs.
\]
\end{theorem}

\begin{proof}
Докажем по индукции по структуре списка $xs$.

\textbf{База.} Пусть $xs = []$.

Тогда:
\[
\operatorname{map}\ f\ (\operatorname{map}\ g\ [])
=
\operatorname{map}\ f\ []
=
[].
\]

С другой стороны:
\[
\operatorname{map}\ (f \circ g)\ []
=
[].
\]

Следовательно, равенство верно.

\textbf{Переход.} Пусть $xs = x:xs'$ и для $xs'$ утверждение верно:
\[
\operatorname{map}\ f\ (\operatorname{map}\ g\ xs')
=
\operatorname{map}\ (f \circ g)\ xs'.
\]

Докажем для $x:xs'$:

\[
\operatorname{map}\ f\ (\operatorname{map}\ g\ (x:xs'))
\]

По определению \texttt{map}:

\[
=
\operatorname{map}\ f\ (g(x) : \operatorname{map}\ g\ xs')
\]

\[
=
f(g(x)) : \operatorname{map}\ f\ (\operatorname{map}\ g\ xs').
\]

По предположению индукции:

\[
=
f(g(x)) : \operatorname{map}\ (f \circ g)\ xs'.
\]

Поскольку
\[
(f \circ g)(x) = f(g(x)),
\]
получаем:

\[
=
(f \circ g)(x) : \operatorname{map}\ (f \circ g)\ xs'
=
\operatorname{map}\ (f \circ g)\ (x:xs').
\]

Следовательно, утверждение доказано.
\end{proof}

\section{Рекурсия и хвостовая рекурсия}

\subsection{Обычная рекурсия}

Функция факториала:

\[
n! =
\begin{cases}
1, & n = 0, \\
n \cdot (n-1)!, & n > 0.
\end{cases}
\]

Псевдокод:

\begin{lstlisting}
fact n =
  if n == 0 then 1
  else n * fact (n - 1)
\end{lstlisting}

Вычисление:

\[
\operatorname{fact}(4)
=
4 \cdot \operatorname{fact}(3)
=
4 \cdot 3 \cdot \operatorname{fact}(2)
=
4 \cdot 3 \cdot 2 \cdot \operatorname{fact}(1)
=
4 \cdot 3 \cdot 2 \cdot 1 \cdot \operatorname{fact}(0)
=
24.
\]

Такая рекурсия требует хранить отложенные умножения.

\subsection{Хвостовая рекурсия}

\begin{definition}
Рекурсивный вызов является \textbf{хвостовым}, если он является последним действием функции.
\end{definition}

Хвостовая версия факториала:

\[
\operatorname{factIter}(n, acc) =
\begin{cases}
acc, & n = 0, \\
\operatorname{factIter}(n-1, acc \cdot n), & n > 0.
\end{cases}
\]

\[
\operatorname{fact}(n) = \operatorname{factIter}(n,1).
\]

Псевдокод:

\begin{lstlisting}
fact n = factIter n 1

factIter n acc =
  if n == 0 then acc
  else factIter (n - 1) (acc * n)
\end{lstlisting}

Иллюстрация вычисления:

\[
\operatorname{factIter}(4,1)
\to
\operatorname{factIter}(3,4)
\to
\operatorname{factIter}(2,12)
\to
\operatorname{factIter}(1,24)
\to
\operatorname{factIter}(0,24)
\to
24.
\]

\subsection{Оптимизация хвостовой рекурсии}

При оптимизации хвостовой рекурсии компилятор может заменить рекурсивный вызов переходом, не создавая новый стековый кадр.

Схематически:

\[
\operatorname{factIter}(n,acc)
\]

заменяется циклом:

\begin{lstlisting}
while n != 0:
    acc = acc * n
    n = n - 1
return acc
\end{lstlisting}

\section{Функциональная декомпозиция}

\subsection{Понятие функциональной декомпозиции}

\begin{definition}
\textbf{Функциональная декомпозиция} --- метод проектирования программы, при котором задача разбивается на набор функций, каждая из которых решает отдельную подзадачу, а итоговое решение получается их композицией.
\end{definition}

Главная идея:

\[
\text{сложная задача}
\quad \Rightarrow \quad
\text{набор простых функций}
\quad \Rightarrow \quad
\text{композиция функций}.
\]

Например, задача обработки списка чисел:

\begin{quote}
Дан список чисел. Нужно оставить только чётные, возвести их в квадрат и сложить.
\end{quote}

Функциональная декомпозиция:

\begin{enumerate}
  \item выбрать чётные числа;
  \item возвести каждое число в квадрат;
  \item сложить элементы списка.
\end{enumerate}

В функциональном стиле:

\[
\operatorname{sum}(\operatorname{map}\ square(\operatorname{filter}\ even\ xs)).
\]

Пример:

\[
xs = [1,2,3,4,5].
\]

Шаги:

\[
[1,2,3,4,5]
\xrightarrow{\operatorname{filter}\ even}
[2,4]
\xrightarrow{\operatorname{map}\ square}
[4,16]
\xrightarrow{\operatorname{sum}}
20.
\]

\subsection{Принципы функциональной декомпозиции}

\begin{enumerate}
  \item \textbf{Разделение вычисления на чистые функции.}

  Каждая функция должна иметь ясную математическую интерпретацию.

  \item \textbf{Минимизация состояния.}

  Состояние по возможности передаётся явно как аргумент и результат.

  \item \textbf{Композиционность.}

  Результат одной функции является входом другой.

  \[
  h = g \circ f.
  \]

  \item \textbf{Обобщение повторяющихся схем.}

  Например, вместо написания рекурсии каждый раз используются \texttt{map}, \texttt{filter}, \texttt{fold}.

  \item \textbf{Изоляция эффектов.}

  Ввод-вывод, работа с файлами, сетью и базой данных отделяются от чистой бизнес-логики.

  \item \textbf{Малые функции.}

  Предпочтительно писать функции, решающие одну понятную задачу.
\end{enumerate}

\subsection{Пример функциональной декомпозиции}

Задача:

\begin{quote}
Для списка заказов вычислить суммарную стоимость оплаченных заказов дороже 1000.
\end{quote}

Пусть заказ имеет поля:

\[
Order = \{id, paid, total\}.
\]

Функции:

\begin{lstlisting}
isPaid order = order.paid

isExpensive order = order.total > 1000

orderTotal order = order.total
\end{lstlisting}

Композиция:

\begin{lstlisting}
result =
  orders
  |> filter isPaid
  |> filter isExpensive
  |> map orderTotal
  |> sum
\end{lstlisting}

Математически:

\[
\operatorname{result}
=
\operatorname{sum}
(
\operatorname{map}\ orderTotal
(
\operatorname{filter}\ isExpensive
(
\operatorname{filter}\ isPaid\ orders
))).
\]

Преимущество такой декомпозиции состоит в том, что каждую функцию можно тестировать независимо.

\subsection{Декомпозиция через конвейер}

В функциональном стиле часто используется конвейерная обработка данных:

\[
x
\to f(x)
\to g(f(x))
\to h(g(f(x))).
\]

Или:

\[
h \circ g \circ f.
\]

Пример:

\[
[1,2,3,4]
\xrightarrow{\operatorname{filter}\ even}
[2,4]
\xrightarrow{\operatorname{map}\ (*10)}
[20,40]
\xrightarrow{\operatorname{fold}\ (+)\ 0}
60.
\]

\subsection{Декомпозиция через алгебраические типы данных}

\begin{definition}
\textbf{Алгебраический тип данных} --- тип, построенный с помощью сумм и произведений типов.
\end{definition}

Тип-произведение соответствует записи:

\[
Person = String \times Int.
\]

Например:

\[
Person(name, age).
\]

Тип-сумма соответствует выбору одного из вариантов:

\[
Shape = Circle(Double) \mid Rectangle(Double, Double).
\]

Пример функции площади:

\[
area(shape) =
\begin{cases}
\pi r^2, & shape = Circle(r), \\
w \cdot h, & shape = Rectangle(w,h).
\end{cases}
\]

Псевдокод:

\begin{lstlisting}
area shape =
  case shape of
    Circle r       -> pi * r * r
    Rectangle w h  -> w * h
\end{lstlisting}

Такой стиль делает структуру данных явной, а обработку --- исчерпывающей.

\section{Неизменяемость и разделение структур}

\subsection{Неизменяемые данные}

\begin{definition}
\textbf{Неизменяемое значение} --- значение, состояние которого нельзя изменить после создания.
\end{definition}

Если требуется ``изменить'' значение, создаётся новое значение.

Например, для списка:

\[
xs = [2,3,4].
\]

Добавление элемента в начало:

\[
ys = 1 : xs = [1,2,3,4].
\]

При этом старый список остаётся доступным:

\[
xs = [2,3,4].
\]

Списки $xs$ и $ys$ разделяют хвост:

\[
ys = 1 \to 2 \to 3 \to 4,
\]

\[
xs = 2 \to 3 \to 4.
\]

\subsection{Структурное разделение}

\begin{definition}
\textbf{Структурное разделение} --- техника, при которой новая версия структуры данных переиспользует неизменённые части старой версии.
\end{definition}

Например, добавление элемента в начало списка требует создать только один новый узел:

\[
ys = Cons(1, xs).
\]

Графически:

\[
ys \to [1] \to [2] \to [3] \to []
\]

\[
xs \quad\quad\ \nearrow [2] \to [3] \to []
\]

Такое разделение делает неизменяемые структуры эффективными.

\section{Персистентные структуры данных}

\subsection{Понятие персистентности}

\begin{definition}
\textbf{Персистентная структура данных} --- структура данных, которая сохраняет доступ к предыдущим версиям после выполнения операций обновления.
\end{definition}

В обычной изменяемой структуре обновление разрушает старое состояние:

\[
S_0 \xrightarrow{update} S_1,
\]

и доступна только версия $S_1$.

В персистентной структуре доступны обе версии:

\[
S_0 \xrightarrow{update} S_1,
\]

\[
S_0 \text{ доступна}, \qquad S_1 \text{ доступна}.
\]

\subsection{Виды персистентности}

Выделяют несколько уровней персистентности.

\begin{enumerate}
  \item \textbf{Частичная персистентность.}

  Можно читать любую старую версию, но обновлять можно только последнюю.

  \[
  S_0 \to S_1 \to S_2 \to S_3.
  \]

  Читать можно $S_0,S_1,S_2,S_3$, но обновлять только $S_3$.

  \item \textbf{Полная персистентность.}

  Можно читать и обновлять любую версию.

  \[
  S_0 \to S_1 \to S_2,
  \]

  и можно породить новую ветвь:

  \[
  S_1 \to S_1'.
  \]

  История становится деревом версий.

  \item \textbf{Конфлюэнтная персистентность.}

  Можно объединять несколько версий в новую.

  Например:

  \[
  S_a, S_b \to S_c.
  \]

  Это требуется, например, для слияния деревьев, графов, документов.

  \item \textbf{Функциональная персистентность.}

  Все операции реализуются без изменения существующих узлов. Старые версии сохраняются автоматически.
\end{enumerate}

\subsection{Персистентный односвязный список}

Список определяется рекурсивно:

\[
List(A) = Nil \mid Cons(A, List(A)).
\]

Операции:

\begin{enumerate}
  \item \texttt{empty} --- пустой список;
  \item \texttt{cons(x, xs)} --- добавить элемент в начало;
  \item \texttt{head(xs)} --- получить первый элемент;
  \item \texttt{tail(xs)} --- получить хвост.
\end{enumerate}

Псевдокод:

\begin{lstlisting}
data List a =
    Nil
  | Cons a (List a)

cons x xs = Cons x xs

head (Cons x xs) = x

tail (Cons x xs) = xs
\end{lstlisting}

Пример:

\[
xs = [2,3],
\]

\[
ys = 1 : xs = [1,2,3].
\]

Старый список:

\[
xs = [2,3]
\]

остаётся доступным.

Оценки сложности:

\[
\operatorname{cons} = O(1),
\]

\[
\operatorname{head} = O(1),
\]

\[
\operatorname{tail} = O(1).
\]

\subsection{Теорема о персистентности списка}

\begin{theorem}
Операция добавления элемента в начало неизменяемого односвязного списка создаёт новую версию списка за $O(1)$ времени и $O(1)$ дополнительной памяти, сохраняя старую версию.
\end{theorem}

\begin{proof}
Пусть имеется список:

\[
xs = Cons(x_1, Cons(x_2, \dots Cons(x_n, Nil)\dots)).
\]

Операция добавления элемента $a$ в начало определяется как:

\[
ys = Cons(a, xs).
\]

При этом создаётся ровно один новый узел $Cons(a, xs)$, содержащий:
\begin{enumerate}
  \item значение $a$;
  \item ссылку на уже существующий список $xs$.
\end{enumerate}

Никакой узел списка $xs$ не изменяется. Следовательно, старая версия $xs$ остаётся доступной.

Поскольку создаётся ровно один узел, расход дополнительной памяти равен $O(1)$. Поскольку операция состоит из создания одного узла и записи одной ссылки, время выполнения равно $O(1)$.

Следовательно, утверждение доказано.
\end{proof}

\subsection{Персистентный стек}

Стек можно реализовать как список.

Операции:

\begin{enumerate}
  \item \texttt{empty};
  \item \texttt{push(x, s)};
  \item \texttt{pop(s)};
  \item \texttt{top(s)}.
\end{enumerate}

Реализация:

\begin{lstlisting}
empty = Nil

push x s = Cons x s

top (Cons x s) = x

pop (Cons x s) = s
\end{lstlisting}

Пример:

\[
s_0 = [],
\]

\[
s_1 = push(1,s_0) = [1],
\]

\[
s_2 = push(2,s_1) = [2,1].
\]

После операции:

\[
pop(s_2) = s_1.
\]

При этом доступны все версии:

\[
s_0 = [], \qquad s_1 = [1], \qquad s_2 = [2,1].
\]

Сложность всех основных операций:

\[
O(1).
\]

\subsection{Персистентное бинарное дерево поиска}

\begin{definition}
\textbf{Бинарное дерево поиска} --- дерево, в котором для каждого узла с ключом $k$:
\begin{itemize}
  \item все ключи в левом поддереве меньше $k$;
  \item все ключи в правом поддереве больше $k$.
\end{itemize}
\end{definition}

Тип дерева:

\[
Tree(A) = Empty \mid Node(Tree(A), A, Tree(A)).
\]

Операция поиска:

\[
member(x, Empty) = false,
\]

\[
member(x, Node(l,y,r)) =
\begin{cases}
true, & x = y, \\
member(x,l), & x < y, \\
member(x,r), & x > y.
\end{cases}
\]

Операция вставки в функциональном стиле:

\[
insert(x, Empty) = Node(Empty,x,Empty),
\]

\[
insert(x, Node(l,y,r)) =
\begin{cases}
Node(l,y,r), & x = y, \\
Node(insert(x,l),y,r), & x < y, \\
Node(l,y,insert(x,r)), & x > y.
\end{cases}
\]

Псевдокод:

\begin{lstlisting}
insert x Empty =
  Node Empty x Empty

insert x (Node l y r) =
  if x == y then
    Node l y r
  else if x < y then
    Node (insert x l) y r
  else
    Node l y (insert x r)
\end{lstlisting}

\subsection{Иллюстрация вставки в персистентное дерево}

Пусть есть дерево:

\[
        5
       / \
      3   8
\]

Вставим $4$.

Путь поиска:

\[
5 \to 3 \to \text{правое поддерево}.
\]

Создаются новые узлы только на пути от корня к месту вставки:

\[
        5'
       /  \
      3'   8
        \
         4
\]

Старые узлы, не лежащие на пути, разделяются между версиями. Узел $8$ используется и старой, и новой версией.

\subsection{Алгоритм вставки в персистентное бинарное дерево поиска}

\textbf{Вход:} ключ $x$, дерево $t$.

\textbf{Выход:} новое дерево $t'$, содержащее $x$.

\begin{enumerate}
  \item Если дерево пусто, создать новый узел с ключом $x$.
  \item Если текущий узел содержит ключ $y$:
  \begin{enumerate}
    \item если $x = y$, вернуть текущий узел;
    \item если $x < y$, рекурсивно вставить $x$ в левое поддерево;
    \item создать новый узел с обновлённым левым поддеревом и старым правым;
    \item если $x > y$, рекурсивно вставить $x$ в правое поддерево;
    \item создать новый узел со старым левым поддеревом и обновлённым правым.
  \end{enumerate}
  \item Вернуть новый корень.
\end{enumerate}

\subsection{Теорема о сложности персистентной вставки в дерево}

\begin{theorem}
Вставка элемента в неизменяемое бинарное дерево поиска высоты $h$ выполняется за $O(h)$ времени и требует $O(h)$ дополнительной памяти.
\end{theorem}

\begin{proof}
Обычный поиск позиции для вставки в бинарном дереве поиска проходит по одному пути от корня к листу. Длина этого пути не превосходит высоты дерева $h$.

В функциональной реализации нельзя изменять существующие узлы. Поэтому для сохранения старой версии требуется создать новые узлы на пути от корня к месту вставки.

Рассмотрим путь:

\[
v_0, v_1, \dots, v_k,
\]

где $v_0$ --- корень, а $v_k$ --- место вставки. Длина пути $k \le h$.

Для каждого узла на этом пути создаётся одна новая копия, в которой одна ссылка указывает на обновлённое поддерево, а другая ссылка переиспользует старое поддерево.

Узлы вне пути не изменяются и разделяются между старой и новой версиями.

Следовательно:
\begin{itemize}
  \item количество посещённых узлов не превосходит $h$;
  \item количество созданных новых узлов не превосходит $h+1$.
\end{itemize}

Значит, время работы равно $O(h)$, а дополнительная память --- $O(h)$.
\end{proof}

\subsection{Персистентные массивы}

Массивы сложнее сделать персистентными, потому что обычное обновление элемента:

\[
a[i] := x
\]

изменяет ячейку на месте.

Наивный функциональный способ:

\[
update(a,i,x) = \text{копия массива } a \text{ с заменой } i\text{-го элемента}.
\]

Сложность:

\[
O(n)
\]

по времени и памяти.

Более эффективный подход --- представить массив как дерево с фиксированным коэффициентом ветвления.

Например, вектор можно представить деревом:

\[
\text{индекс} \to \text{путь в дереве}.
\]

При обновлении копируется только путь от корня к листу.

Если коэффициент ветвления равен $b$, а длина массива $n$, высота дерева:

\[
h = O(\log_b n).
\]

Тогда обновление требует:

\[
O(\log_b n)
\]

времени и памяти.

\subsection{Персистентный вектор как дерево}

Идея:

\begin{enumerate}
  \item Разбить индекс на группы битов.
  \item Каждая группа выбирает переход на очередном уровне дерева.
  \item Лист содержит значение.
  \item При обновлении копируется путь от корня до листа.
  \item Неизменённые ветви разделяются.
\end{enumerate}

Пример с ветвлением $b = 4$:

\[
n = 16,\qquad h = \log_4 16 = 2.
\]

Индекс $i$ представляется двумя цифрами в системе счисления по основанию $4$:

\[
i = d_1d_0.
\]

Первая цифра выбирает дочерний узел корня, вторая --- позицию в листе.

\section{Подходы к проектированию функциональных программ}

\subsection{Проектирование от типов}

\begin{definition}
\textbf{Проектирование от типов} --- подход, при котором сначала описываются типы предметной области, затем функции над ними.
\end{definition}

Алгоритм проектирования:

\begin{enumerate}
  \item Определить основные сущности предметной области.
  \item Описать их типами.
  \item Описать допустимые состояния явно.
  \item Исключить невозможные состояния с помощью системы типов.
  \item Реализовать функции, соответствующие переходам между состояниями.
\end{enumerate}

Пример.

Плохо:

\begin{lstlisting}
status : String
\end{lstlisting}

Такой тип допускает любые строки:

\[
"paid",\quad "cancelled",\quad "hello",\quad "".
\]

Лучше:

\begin{lstlisting}
data OrderStatus =
    Created
  | Paid
  | Cancelled
  | Shipped
\end{lstlisting}

Теперь невозможные состояния не выражаются.

\subsection{Разделение чистого ядра и эффектной оболочки}

Практическая функциональная архитектура часто строится следующим образом:

\[
\text{ввод} \to \text{чистое ядро} \to \text{вывод}.
\]

\begin{enumerate}
  \item \textbf{Эффектная оболочка} читает данные из внешнего мира.
  \item \textbf{Чистое ядро} выполняет вычисления.
  \item \textbf{Эффектная оболочка} записывает результат.
\end{enumerate}

Пример:

\begin{lstlisting}
main =
  readFile "orders.json"
  |> parseOrders
  |> calculateReport
  |> writeFile "report.json"
\end{lstlisting}

Здесь функции \texttt{parseOrders} и \texttt{calculateReport} желательно делать чистыми.

\subsection{Функциональное ядро, императивная оболочка}

Этот архитектурный принцип особенно полезен в промышленной разработке.

\[
\boxed{\text{Functional Core}} + \boxed{\text{Imperative Shell}}
\]

\textbf{Функциональное ядро}:
\begin{itemize}
  \item содержит бизнес-логику;
  \item не имеет побочных эффектов;
  \item легко тестируется;
  \item оперирует неизменяемыми значениями.
\end{itemize}

\textbf{Императивная оболочка}:
\begin{itemize}
  \item работает с сетью, базой данных, файлами;
  \item вызывает функции ядра;
  \item выполняет побочные эффекты.
\end{itemize}

\subsection{Декларативное описание вычислений}

В функциональном программировании предпочтительно описывать результат, а не последовательность мутаций.

Императивно:

\begin{lstlisting}
result = []
for x in xs:
    if x % 2 == 0:
        result.append(x * x)
\end{lstlisting}

Функционально:

\begin{lstlisting}
result = map square (filter even xs)
\end{lstlisting}

Или в виде композиции:

\[
result = \operatorname{map}\ square \circ \operatorname{filter}\ even.
\]

\subsection{Использование рекурсивных схем}

Многие программы над структурами данных имеют одинаковую форму. Эти формы можно выделить.

\subsubsection{Катаморфизм}

\begin{definition}
\textbf{Катаморфизм} --- обобщённая свёртка, разбирающая рекурсивную структуру данных.
\end{definition}

Для списка катаморфизмом является \texttt{foldr}.

Например, сумма списка:

\[
sum = foldr\ (+)\ 0.
\]

Длина списка:

\[
length = foldr\ (\lambda x\ acc. acc + 1)\ 0.
\]

Проверка существования элемента:

\[
any\ p = foldr\ (\lambda x\ acc. p(x) \lor acc)\ false.
\]

\subsubsection{Анаморфизм}

\begin{definition}
\textbf{Анаморфизм} --- обобщённое разворачивание структуры из начального состояния.
\end{definition}

Например, генерация списка чисел от $n$ до $0$:

\[
unfold(n) =
\begin{cases}
[], & n < 0, \\
n : unfold(n-1), & n \ge 0.
\end{cases}
\]

\subsubsection{Гиломорфизм}

\begin{definition}
\textbf{Гиломорфизм} --- композиция анаморфизма и катаморфизма: сначала структура порождается, затем сворачивается.
\end{definition}

Например, факториал можно понимать как:

\[
n \to [1,2,\dots,n] \to \operatorname{product}.
\]

\[
factorial(n) = \operatorname{foldr}\ (*)\ 1\ [1,2,\dots,n].
\]

\subsection{Ленивые вычисления}

\begin{definition}
\textbf{Ленивые вычисления} --- стратегия вычислений, при которой выражение вычисляется только тогда, когда его значение действительно требуется.
\end{definition}

Пример бесконечного списка:

\[
ones = 1 : ones.
\]

В строгом языке такое определение привело бы к бесконечному вычислению. В ленивом языке можно взять конечный префикс:

\[
take(5, ones) = [1,1,1,1,1].
\]

\subsection{Строгая и ленивая стратегии вычислений}

\begin{definition}
\textbf{Строгая стратегия} вычисляет аргументы функции до вызова функции.
\end{definition}

\begin{definition}
\textbf{Ленивая стратегия} передаёт аргументы как отложенные вычисления и вычисляет их только при необходимости.
\end{definition}

Пример:

\[
const(x,y) = x.
\]

Рассмотрим:

\[
const(1, \Omega),
\]

где

\[
\Omega = (\lambda x. x x)(\lambda x. x x)
\]

--- расходящийся терм.

При строгой стратегии вычисление $\Omega$ не завершится.

При ленивой стратегии:

\[
const(1,\Omega) \to 1.
\]

\subsection{Мемоизация}

\begin{definition}
\textbf{Мемоизация} --- сохранение результатов вычислений для повторного использования.
\end{definition}

В функциональном программировании мемоизация особенно естественна, потому что чистые функции при одинаковых аргументах всегда возвращают одинаковый результат.

Пример чисел Фибоначчи:

\[
F(0)=0,\quad F(1)=1,
\]

\[
F(n)=F(n-1)+F(n-2).
\]

Наивная рекурсия:

\begin{lstlisting}
fib n =
  if n <= 1 then n
  else fib (n - 1) + fib (n - 2)
\end{lstlisting}

Повторно вычисляет одни и те же значения.

Алгоритм с мемоизацией:

\begin{enumerate}
  \item Создать таблицу уже вычисленных значений.
  \item При вызове \texttt{fib(n)} проверить, есть ли результат в таблице.
  \item Если есть, вернуть его.
  \item Если нет, вычислить рекурсивно.
  \item Сохранить результат в таблицу.
  \item Вернуть результат.
\end{enumerate}

Для чистой функции можно рассматривать таблицу как неявный кэш, не меняющий математический смысл функции.

\section{Реализация функциональных программ}

\subsection{Замыкания}

\begin{definition}
\textbf{Замыкание} --- функция вместе с окружением, в котором она была создана.
\end{definition}

Пример:

\begin{lstlisting}
makeAdder a =
  let f x = x + a
  in f
\end{lstlisting}

Вызов:

\[
add5 = makeAdder(5),
\]

\[
add5(3) = 8.
\]

Функция \texttt{add5} хранит не только код \texttt{f}, но и значение $a = 5$.

Схематически:

\[
closure = \langle code,\ environment \rangle.
\]

\subsection{Реализация функций высшего порядка}

Функция высшего порядка принимает замыкание как обычное значение.

Например:

\[
map(f, xs).
\]

Фактически \texttt{f} передаётся как структура:

\[
\langle code_f,\ environment_f \rangle.
\]

При применении:

\[
f(x)
\]

используется:
\begin{enumerate}
  \item код функции;
  \item сохранённое окружение;
  \item текущий аргумент.
\end{enumerate}

\subsection{Сборка мусора}

Неизменяемые и персистентные структуры часто создают много краткоживущих объектов. Поэтому функциональные языки обычно полагаются на автоматическое управление памятью.

\begin{definition}
\textbf{Сборка мусора} --- автоматическое освобождение памяти, занятой объектами, недостижимыми из программы.
\end{definition}

Для персистентных структур важно, что узел можно удалить только тогда, когда на него не ссылается ни одна актуальная версия.

\subsection{Оптимизация хвостовых вызовов}

\begin{definition}
\textbf{Оптимизация хвостовых вызовов} --- замена вызова функции в хвостовой позиции переходом без роста стека вызовов.
\end{definition}

Например:

\[
f(x) =
\begin{cases}
x, & x = 0, \\
f(x-1), & x > 0.
\end{cases}
\]

Рекурсивный вызов $f(x-1)$ находится в хвостовой позиции.

\subsection{Каррирование}

\begin{definition}
\textbf{Каррирование} --- преобразование функции от нескольких аргументов в цепочку функций от одного аргумента.
\end{definition}

Функция:

\[
f : A \times B \to C
\]

заменяется функцией:

\[
f' : A \to (B \to C).
\]

Пример:

\[
add(x,y)=x+y.
\]

Каррированная версия:

\[
add(x)(y)=x+y.
\]

Тогда:

\[
add(5)
\]

--- функция, прибавляющая $5$.

\subsection{Частичное применение}

\begin{definition}
\textbf{Частичное применение} --- создание новой функции путём фиксации части аргументов исходной функции.
\end{definition}

Пример:

\[
add5 = add(5).
\]

Тогда:

\[
add5(10) = 15.
\]

\section{Монады и управление эффектами}

\subsection{Проблема эффектов}

Чистые функции не должны выполнять ввод-вывод, генерировать случайные числа, изменять глобальное состояние. Однако реальные программы взаимодействуют с внешним миром.

Один из подходов --- представить эффектное вычисление как значение специального типа.

Например:

\[
IO\ A
\]

можно понимать как описание действия, которое при выполнении может взаимодействовать с внешним миром и вернуть значение типа $A$.

\subsection{Понятие монады}

\begin{definition}
\textbf{Монада} --- типовой конструктор $M$ вместе с операциями:
\[
return : A \to M A,
\]
\[
bind : M A \to (A \to M B) \to M B,
\]
удовлетворяющими монадическим законам.
\end{definition}

Обычно \texttt{bind} записывается как:

\[
>>=.
\]

Интуитивно:
\begin{itemize}
  \item \texttt{return} помещает чистое значение в контекст;
  \item \texttt{bind} связывает вычисление в контексте со следующим шагом.
\end{itemize}

\subsection{Законы монад}

Для монады должны выполняться три закона.

\subsubsection{Левая единица}

\[
return\ a >>= f = f(a).
\]

\subsubsection{Правая единица}

\[
m >>= return = m.
\]

\subsubsection{Ассоциативность}

\[
(m >>= f) >>= g
=
m >>= (\lambda x. f(x) >>= g).
\]

Эти законы гарантируют согласованное поведение последовательной композиции эффектных вычислений.

\subsection{Пример: Maybe}

Тип:

\[
Maybe\ A = Nothing \mid Just(A).
\]

Он моделирует вычисления, которые могут завершиться неудачей.

Операция \texttt{bind}:

\[
Nothing >>= f = Nothing,
\]

\[
Just(x) >>= f = f(x).
\]

Пример безопасного деления:

\[
safeDiv(x,y) =
\begin{cases}
Nothing, & y = 0, \\
Just(x/y), & y \ne 0.
\end{cases}
\]

Композиция:

\[
Just(10) >>= (\lambda x. safeDiv(x,2)) >>= (\lambda y. safeDiv(y,5)).
\]

Шаги:

\[
Just(10) >>= safeDiv(\cdot,2) = Just(5),
\]

\[
Just(5) >>= safeDiv(\cdot,5) = Just(1).
\]

Если возникает деление на ноль:

\[
Just(10) >>= safeDiv(\cdot,0) = Nothing,
\]

и дальнейшие вычисления не выполняются.

\section{Подходы к реализации персистентных структур данных}

\subsection{Копирование пути}

\begin{definition}
\textbf{Копирование пути} --- способ реализации персистентного обновления дерева, при котором копируются только узлы на пути к изменяемому месту.
\end{definition}

Алгоритм:

\begin{enumerate}
  \item Найти путь от корня к изменяемому узлу.
  \item Создать новую копию изменяемого узла.
  \item Поднимаясь вверх, создать копии всех предков.
  \item Неизменённые поддеревья переиспользовать.
  \item Вернуть новый корень.
\end{enumerate}

Пример:

\[
        A
       / \
      B   C
     / \
    D   E
\]

Изменим $E$ на $E'$.

Новая версия:

\[
        A'
       /  \
      B'   C
     /  \
    D    E'
\]

Узлы $C$ и $D$ разделяются между версиями.

\subsection{Жирные узлы}

\begin{definition}
\textbf{Жирные узлы} --- способ реализации частично персистентных структур, при котором в каждом узле хранятся не только текущие значения полей, но и история их изменений.
\end{definition}

Например, поле \texttt{left} может храниться как список пар:

\[
[(version_1, value_1), (version_2, value_2), \dots].
\]

Чтение поля в версии $v$ означает поиск последнего изменения с номером версии не больше $v$.

Алгоритм чтения поля:

\begin{enumerate}
  \item Взять список изменений поля.
  \item Найти изменение с максимальной версией $u \le v$.
  \item Вернуть соответствующее значение.
\end{enumerate}

Преимущество:
\begin{itemize}
  \item не требуется копировать весь путь при каждом изменении.
\end{itemize}

Недостаток:
\begin{itemize}
  \item чтение полей может стать дороже;
  \item узлы усложняются;
  \item требуется хранить версии.
\end{itemize}

\subsection{Метод копирования узлов}

Метод копирования узлов используется для ограниченной степени входа в указательных структурах.

Идея:

\begin{enumerate}
  \item Каждый узел имеет ограниченное число дополнительных ячеек для изменений.
  \item Если есть свободная ячейка, изменение записывается в неё.
  \item Если свободных ячеек нет, узел копируется.
  \item Изменение родительского указателя распространяется вверх.
\end{enumerate}

Такой подход позволяет эффективно делать указательные структуры частично или полностью персистентными при ограничениях на степень узлов.

\section{Сравнение функционального и императивного подходов}

\subsection{Работа с состоянием}

В императивном стиле состояние изменяется явно:

\begin{lstlisting}
x = x + 1
\end{lstlisting}

В функциональном стиле создаётся новое значение:

\begin{lstlisting}
x2 = x + 1
\end{lstlisting}

\subsection{Циклы и рекурсия}

Императивно:

\begin{lstlisting}
s = 0
for x in xs:
    s += x
\end{lstlisting}

Функционально:

\begin{lstlisting}
sum xs = foldr (+) 0 xs
\end{lstlisting}

\subsection{Обновление структур данных}

Императивное дерево:

\[
T \xrightarrow{insert} T
\]

То же дерево изменяется на месте.

Функциональное дерево:

\[
T \xrightarrow{insert} T'
\]

Старая версия $T$ остаётся доступной.

\subsection{Тестирование}

Чистая функция:

\[
f : A \to B
\]

тестируется как соответствие входа и выхода.

Нечистая функция зависит от:

\begin{itemize}
  \item глобальных переменных;
  \item времени;
  \item файлов;
  \item сети;
  \item состояния базы данных;
  \item порядка вызовов.
\end{itemize}

Поэтому её тестирование сложнее.

\section{Минимальные примеры функциональных программ}

\subsection{Сумма квадратов чётных чисел}

Задача:

\[
xs = [1,2,3,4,5,6].
\]

Найти:

\[
2^2 + 4^2 + 6^2 = 56.
\]

Функциональное решение:

\begin{lstlisting}
sumSquaresEven xs =
  xs
  |> filter even
  |> map square
  |> sum
\end{lstlisting}

Шаги:

\[
[1,2,3,4,5,6]
\to
[2,4,6]
\to
[4,16,36]
\to
56.
\]

\subsection{Подсчёт слов}

Задача: по строке получить частоты слов.

Функциональная декомпозиция:

\begin{enumerate}
  \item привести строку к нижнему регистру;
  \item удалить пунктуацию;
  \item разбить на слова;
  \item сгруппировать одинаковые слова;
  \item посчитать элементы в группах.
\end{enumerate}

Композиция:

\begin{lstlisting}
wordCount text =
  text
  |> toLower
  |> removePunctuation
  |> splitWords
  |> group
  |> map countGroup
\end{lstlisting}

\subsection{Безопасная обработка данных}

Пусть нужно:
\begin{enumerate}
  \item прочитать строку;
  \item преобразовать её в число;
  \item взять обратное значение.
\end{enumerate}

Проблемы:
\begin{itemize}
  \item строка может не быть числом;
  \item число может быть нулём.
\end{itemize}

Используем \texttt{Maybe}:

\begin{lstlisting}
parseNumber : String -> Maybe Number
reciprocal  : Number -> Maybe Number

process s =
  parseNumber s >>= reciprocal
\end{lstlisting}

Примеры:

\[
process("4") = Just(0.25),
\]

\[
process("0") = Nothing,
\]

\[
process("abc") = Nothing.
\]

\section{Общие рекомендации по проектированию функциональных программ}

\begin{enumerate}
  \item Сначала проектировать данные и типы.
  \item Делать невозможные состояния непредставимыми.
  \item Выносить побочные эффекты на границы программы.
  \item Бизнес-логику реализовывать чистыми функциями.
  \item Использовать композицию вместо длинных процедур.
  \item Повторяющиеся рекурсивные схемы заменять на \texttt{map}, \texttt{filter}, \texttt{fold}.
  \item Предпочитать неизменяемые структуры данных.
  \item Для эффективности использовать структурное разделение.
  \item Для больших структур использовать деревья, векторы с ветвлением, сбалансированные деревья.
  \item Анализировать сложность с учётом количества создаваемых новых узлов.
  \item Использовать хвостовую рекурсию там, где возможна линейная обработка.
  \item Отделять описание вычисления от его интерпретации.
\end{enumerate}

\section{Краткие выводы}

Функциональное программирование основывается на представлении программы как композиции функций. Главными понятиями являются чистые функции, неизменяемость, ссылочная прозрачность, функции высшего порядка, рекурсия и композиция.

Функциональная декомпозиция позволяет строить программы из малых независимых функций. Это повышает модульность, тестируемость и понятность кода.

Персистентные структуры данных сохраняют доступ к старым версиям после обновлений. Их эффективность достигается за счёт структурного разделения: новые версии переиспользуют неизменённые части старых структур. Для списков и стеков многие операции выполняются за $O(1)$, для деревьев обновление обычно требует $O(h)$ времени и памяти, где $h$ --- высота дерева.

Проектирование функциональных программ обычно строится вокруг типов, чистого ядра, изоляции эффектов, функций высшего порядка и рекурсивных схем обработки данных.

\end{document}
